\documentclass[12pt,onecolumn]{article}

\PassOptionsToPackage{capitalise,nameinlink,noabbrev}{cleveref}

\usepackage[utf8]{inputenc}
\usepackage[T1]{fontenc}
\usepackage{lmodern}
\usepackage[margin=1in]{geometry}
\usepackage{setspace}
\usepackage{microtype}
\usepackage{amsmath,amssymb,amsthm,bm}
\usepackage{booktabs}
\usepackage[numbers,sort&compress]{natbib}
\usepackage[colorlinks=true,linkcolor=blue,citecolor=blue,urlcolor=blue]{hyperref}

\newcommand{\doi}[1]{\href{https://doi.org/#1}{#1}}

\newtheorem{proposition}{Proposition}[section]
\newtheorem{definition}[proposition]{Definition}
\newtheorem{lemma}[proposition]{Lemma}
\newtheorem{corollary}[proposition]{Corollary}
\newtheorem{remark}[proposition]{Remark}

\usepackage{cleveref}
\usepackage{orcidlink}
\usepackage{fancyhdr}

\newcommand\blfootnote[1]{%
  \begingroup
  \renewcommand\thefootnote{}\footnote{#1}%
  \addtocounter{footnote}{-1}%
  \endgroup
}

\setstretch{1.08}

\pagestyle{fancy}
\fancyhf{}
\setlength{\headheight}{14pt}
\lhead{The Present Was Never Universal}
\rhead{John C. W. McKinley}
\cfoot{\thepage}

\title{\textbf{The Present Was Never Universal}\\
\large Local Registration and the Absence of a Universe-Wide Now}

\author{John C. W. McKinley\,\orcidlink{0009-0005-7097-5035}}

\date{July 2026}

\begin{document}

\maketitle

\blfootnote{\scriptsize Draft version. DOI placeholder to be supplied upon publication.}

\begin{abstract}
This paper makes one narrow interpretive claim: the present is not a universal physical condition shared by the universe as a whole. The present is unique at every indivisible local site in the universe. Each such micro-site occupies its own spacetime location, has its own causal position, and admits only the information and lawful physical relations available there. A detector or observer therefore does not originate locality. It integrates an immense plurality of micro-local presents into a higher-order registration presented as a world. Other detectors do not receive the same present from different viewpoints. They integrate different collections of micro-local presents whose relations are constrained to remain mutually consistent by special relativity, general relativity, and the invariant causal structure conventionally expressed by \(c\). On the Timeless Light Model reading, the information speed limit does more than prevent superluminal transmission. It limits which portions of the universe can be presented together to a detector and thereby produces the ordered appearance of distance, depth, delay, persistence, and shared history. A macroscopic body appears continuously present because enormous numbers of local registrations preserve a stable relational pattern. At smaller scales, the same body is represented through probability-governed lawful outcomes rather than through an infinitely precise, universally present material state. The result is not that the world is unreal. The result is that physical presence is registration-local. The past was never a universal present that later disappeared, and the present is never a universe-wide state that exists before registration. The present is unique at every indivisible micro-local site, while an observer or detector integrates many such local presents into a registration presented as a world. The claim is interpretive only and alters no equations or predictions of special relativity, general relativity, quantum mechanics, or quantum field theory.
\end{abstract}

\noindent\textbf{Keywords:} present; simultaneity; observer frame; detector; registration; special relativity; general relativity; information speed limit; quantum probability; Timeless Light Model.

\section{Introduction}

Ordinary language treats the present as though it were a universal condition. The universe is imagined to possess one complete state called ``now,'' and observers are imagined to occupy different locations within that same shared moment. Under this picture, the past is the former universal present, the present is the current universal state, and the future is the universal state that has not yet arrived.

Relativity does not supply that architecture. Special relativity denies a frame-independent simultaneity relation for spatially separated events~\cite{Einstein1905,Rindler2006}. General relativity further restricts temporal description by making clock rates and simultaneity assignments dependent on spacetime geometry, motion, and the chosen foliation~\cite{Wald1984,MTW1973}. A detector does not receive the universe as it exists ``all at once.'' It receives only those causal influences that reach its spacetime location. Its physical record is local.

The Timeless Light Model has already stated the corresponding observer-side interpretation of the information speed limit. The constant \(c\) is not a traveler-owned speed but the invariant causal bound governing how lightlike relations, events, and cross-frame relations are registered in spacetime description~\cite{RegistrationBound}. The observer-perceivability paper further identifies \(c\) as the structural bound on the spacetime registration of lawful state changes~\cite{ObserverPerceivability}. Once those results are accepted, the status of the present should be stated explicitly.

The present paper supplies that statement:

\begin{quote}
\textbf{The present was never a universal physical condition. It is a local act of registration presented to one observer or detector as a world. Every other observer or detector registers a different present.}
\end{quote}

The phrase ``different present'' does not mean that physical law becomes private or arbitrary. Local difference begins below the detector. Every indivisible micro-local site occupies its own causal position and therefore has its own present. A detector occupies and coordinates an extended collection of such sites, producing one higher-order registration from many micro-local presents. The registrations remain lawfully related. Relativity supplies the transformation rules, causal structure, geometric relations, and consistency conditions by which the different local presents can be compared.

The paper distinguishes four claims.

\begin{enumerate}
\item There is no universe-wide physical present.
\item The present is unique at every indivisible micro-local site.
\item A detector present is an organized integration of many micro-local presents.
\item Different detectors integrate different local presents.
\item The apparent shared world is the lawful consistency among those registrations, not an independently given universal now.
\end{enumerate}

The argument is interpretive. It introduces no new equation, modifies no tested prediction, and requires no departure from standard relativistic or quantum formalism. Its contribution is to remove the residual ontology of a universal present from a framework in which neither relativity nor detector physics licenses one.

\section{Prior Results and Scope}

The present paper depends on a narrow set of prior results.

\begin{proposition}[Prior registration results]
\label{prop:prior-results}
The following results are presupposed:

\begin{enumerate}
\item a photon is not a persistent traveler carrying an internal history through spacetime, but a lawfully admissible charge-state relation whose spacetime appearance is a lawful change~\cite{Bedrock};
\item wavefunction prediction does not license an intermediate photon path~\cite{Timebound};
\item \(c\) is the invariant causal bound governing observer-side registration of cross-frame and lightlike relations, not a traveler-owned speed~\cite{RegistrationBound};
\item observer-perceivability is spacetime-side registration at a detector, instrument, record-bearing interaction, or observer~\cite{ObserverPerceivability}.
\end{enumerate}
\end{proposition}

\begin{remark}
The present paper does not re-argue photon ontology. It begins after the traveler reading has already been removed. The question is what follows for the temporal status of the world presented at registration.
\end{remark}

The scope is equally narrow. The paper does not deny the usefulness of coordinate time, simultaneity conventions, spacelike hypersurfaces, cosmological time parameters, or frame-dependent descriptions of extended systems. It denies that any such mathematical construction constitutes a single observer-independent physical present already possessed by the universe.

\section{Definitions}

\begin{definition}[Registration event]
\label{def:registration}
A registration event is a physically definite spacetime event at which a detector, observer, instrument, material system, or record-bearing interaction acquires a determinate local outcome or record.
\end{definition}

\begin{definition}[Detector]
\label{def:detector}
A detector is any physical system capable of participating in a registration event. The term includes laboratory instruments, photographic plates, antennas, biological sensory systems, interacting atoms, and other systems in which a definite physical record is produced.
\end{definition}

\begin{definition}[Micro-local site]
\label{def:micro-site}
A micro-local site is the smallest physically indivisible location at which a distinct local physical state, causal relation, or registration possibility can be instantiated. It is the limiting unit of physical locality in the present argument. The paper does not identify that unit with any particular existing theory of quantum gravity.
\end{definition}

\begin{definition}[Micro-local present]
\label{def:micro-present}
The micro-local present is the unique lawful physical condition instantiated at one micro-local site, including only the causal information, field relations, interaction possibilities, and registered state available at that site.
\end{definition}

\begin{definition}[Integrated detector present]
\label{def:integrated-present}
An integrated detector present is the organized higher-order registration produced by the coordinated interaction of the many micro-local sites composing a detector or observer.
\end{definition}



\begin{definition}[Local present]
\label{def:local-present}
The local present of a detector is its integrated detector present: the determinate higher-order registration produced across the micro-local sites composing the detector, together with the locally available physical records and causal information instantiated through that organized system.
\end{definition}

\begin{definition}[Universal present]
\label{def:universal-present}
A universal present is a purported observer-independent physical condition in which all spatially separated regions of the universe possess one jointly existing and uniquely defined state of ``now.''
\end{definition}

\begin{definition}[Registration world]
\label{def:registration-world}
A registration world is the coherent world presented to a detector through its current registration, its retained records, and the lawful relations among the events accessible from its spacetime location.
\end{definition}

\begin{definition}[Shared world]
\label{def:shared-world}
A shared world is the network of different local registration worlds whose records, causal relations, and coordinate descriptions remain mutually transformable and lawfully consistent.
\end{definition}

\section{Relativity Withholds a Universal Present}

\begin{proposition}[No frame-independent distant simultaneity]
\label{prop:no-global-simultaneity}
Special relativity does not provide a unique frame-independent relation of simultaneity among spatially separated events.
\end{proposition}

\begin{proof}
Let two spatially separated events have coordinate differences \(\Delta t\) and \(\Delta x\) in an inertial frame \(S\). Under a Lorentz transformation to an inertial frame \(S'\) moving with velocity \(v\) along the \(x\)-axis,
\begin{equation}
\Delta t'=\gamma\left(\Delta t-\frac{v\Delta x}{c^2}\right),
\qquad
\gamma=\frac{1}{\sqrt{1-v^2/c^2}}.
\end{equation}
If the events are simultaneous in \(S\), then \(\Delta t=0\), but
\begin{equation}
\Delta t'=-\gamma\frac{v\Delta x}{c^2},
\end{equation}
which is nonzero whenever \(v\neq 0\) and \(\Delta x\neq 0\). Therefore spatially separated simultaneity in one inertial frame is not simultaneity in every inertial frame. No unique inertial-frame-independent distant present is supplied by the formalism.
\end{proof}

\begin{corollary}[The universal present is not a relativistic invariant]
\label{cor:not-invariant}
The universal present of \cref{def:universal-present} is not a relativistic invariant.
\end{corollary}

\begin{proof}
A universal present requires a unique observer-independent assignment of ``now'' to spatially separated regions. By \cref{prop:no-global-simultaneity}, special relativity supplies no such assignment. Therefore the universal present is not a relativistic invariant.
\end{proof}

\begin{proposition}[General relativity preserves locality]
\label{prop:gr-locality}
General relativity does not restore a unique universe-wide physical present.
\end{proposition}

\begin{proof}
General relativity describes spacetime through a differentiable manifold equipped with a metric tensor \(g_{\mu\nu}\). Local causal structure is defined by the metric light cones. Extended assignments of simultaneity require coordinate choices or foliations of spacetime into spacelike hypersurfaces. Different admissible coordinate systems and foliations assign different sets of distant events to the same coordinate time. In generic curved spacetimes, no preferred global foliation is supplied by the field equations alone~\cite{Wald1984,Malament2012}. Therefore general relativity preserves local causal and temporal structure without providing a unique observer-independent physical present across the universe.
\end{proof}

\begin{remark}
Some cosmological models admit a useful cosmic-time parameter tied to a preferred family of comoving observers. That parameter supplies a powerful large-scale description. It does not convert all spatially separated events into one directly registered physical present. Each actual detector still occupies one local spacetime event and receives only its own causal information.
\end{remark}

\section{The Present Belongs to Registration}

\begin{proposition}[Registration locality]
\label{prop:registration-locality}
Every actual physical registration occurs locally at one spacetime event or within one bounded local interaction region.
\end{proposition}

\begin{proposition}[Micro-local uniqueness]
\label{prop:micro-uniqueness}
Every micro-local site has a unique present.
\end{proposition}

\begin{proof}
A micro-local site occupies a distinct physical location. Its causal past, local field relations, interaction opportunities, and immediately instantiated physical state are therefore specific to that site. No numerically distinct site occupies precisely the same physical location or receives precisely the same set of local relations. By \cref{def:micro-present}, the lawful condition instantiated there is its micro-local present. Therefore every micro-local site has a unique present.
\end{proof}

\begin{proposition}[Detector integration]
\label{prop:detector-integration}
A detector present is an organized integration of many micro-local presents.
\end{proposition}

\begin{proof}
A detector is an extended physical system composed of many locally situated constituents. Each constituent site has its own micro-local present by \cref{prop:micro-uniqueness}. Through lawful coupling, signal transfer, material constraint, and record formation, those sites contribute to one coordinated detector response. By \cref{def:integrated-present}, that coordinated response is the detector's integrated present. Therefore a detector present is an organized integration of many micro-local presents.
\end{proof}

\begin{proof}
A physical record requires a definite interaction: absorption at a sensor, charge displacement in a circuit, chemical change on a plate, excitation of an atom, a memory-state change, or another localized physical response. Such interactions occur where the participating systems meet in spacetime. No detector records an event from nowhere, and no detector directly records the entire universe at one instant. Therefore every actual physical registration is local.
\end{proof}

\begin{proposition}[The present is registration-local]
\label{prop:present-local}
A physical present exists first as the unique present of a micro-local site and, at detector scale, as the organized integration of many such presents.
\end{proposition}

\begin{proof}
By \cref{cor:not-invariant,prop:gr-locality}, relativity supplies no unique universe-wide physical present. By \cref{prop:registration-locality}, every actual physical record occurs locally. What is physically definite begins with the micro-local condition instantiated at each site. By \cref{prop:micro-uniqueness}, each site has its own present. By \cref{prop:detector-integration}, a detector organizes many such presents into one higher-order registration. No additional universe-wide present is required or physically registered. Therefore physical presence is micro-locally unique and detector-level presence is integrated rather than primitive.
\end{proof}

\begin{corollary}[The present is not prior to registration]
\label{cor:not-prior}
The present is not a complete universe-wide condition that exists first and is then observed locally. The physically available present is constituted at local registration.
\end{corollary}

\begin{proof}
If a complete universal present existed prior to registration, it would require the observer-independent distant simultaneity excluded by \cref{cor:not-invariant}. What is physically instantiated at the detector is the local registration described in \cref{prop:present-local}. Therefore the present is not a pre-existing universe-wide state later sampled by the detector. It is the local registered condition itself.
\end{proof}

\section{Every Detector Registers a Different Present}

\begin{proposition}[Distinct detector events]
\label{prop:distinct-events}
Two spatially separated detectors do not occupy the same registration event.
\end{proposition}

\begin{proof}
A spacetime event is specified by its location in spacetime. Spatially separated detectors occupy different spatial locations and therefore different events, even when one coordinate system assigns equal coordinate time to their registrations. Hence they do not occupy the same registration event.
\end{proof}

\begin{proposition}[Different causal pasts]
\label{prop:different-pasts}
Spatially separated detectors possess different causal pasts and therefore receive different sets of lawfully available information.
\end{proposition}

\begin{proof}
The causal past of an event is the set of events capable of influencing it by timelike or null relations. Two distinct detector events have distinct past light cones. The cones overlap but are not identical. Each detector therefore receives some information unavailable at the other detector's event, and each lacks some information available to the other. Their locally instantiated records are not identical.
\end{proof}

\begin{proposition}[Detector-relative present]
\label{prop:detector-relative}
Every detector integrates a different higher-order present from its own plurality of micro-local presents.
\end{proposition}

\begin{proof}
Different detectors are composed of different collections of micro-local sites. By \cref{prop:micro-uniqueness}, each of those sites has its own present. By \cref{prop:distinct-events,prop:different-pasts}, the detectors also occupy different causal positions and receive different causal information. Their motion, orientation, gravitational environment, detector coupling, and prior records further determine how their constituent micro-local presents are integrated. Therefore each detector produces a different higher-order present.
\end{proof}

\begin{remark}
``Different'' does not mean contradictory. Detector \(A\) and detector \(B\) need not produce the same local record in order to inhabit one lawful universe. Their records must instead be related by the applicable causal, geometric, and dynamical laws.
\end{remark}

\begin{corollary}[No view from nowhere]
\label{cor:no-nowhere}
There is no physically registered view from nowhere in which all detector-relative presents are replaced by one absolute present.
\end{corollary}

\begin{proof}
Every registration is local by \cref{prop:registration-locality}, and every detector registers from a distinct event by \cref{prop:detector-relative}. A view from nowhere would be a registration without a registering event, causal position, or frame. No such physical registration exists.
\end{proof}

\section{The Information Speed Limit Creates Presented Depth}

The information speed limit has usually been described negatively: no influence or information propagates faster than \(c\). The registration reading supplies a positive architectural role. The bound determines which distant relations can enter a local present and in what causal order.

\begin{definition}[Presented depth]
\label{def:presented-depth}
Presented depth is the spatial and temporal ordering experienced or recorded by a detector because information from differently located events becomes available only according to the causal structure bounded by \(c\).
\end{definition}

\begin{proposition}[Finite causal access]
\label{prop:finite-access}
At any registration event, a detector has access only to information contained within its causal past.
\end{proposition}

\begin{proof}
An event outside the detector event's past light cone is spacelike separated from it and cannot have transmitted a causal influence to it without violating the spacetime causal bound. Therefore such an event cannot yet appear in the detector's local registration. The detector's physically available information is restricted to its causal past.
\end{proof}

\begin{proposition}[Depth from bounded registration]
\label{prop:depth}
The ordered appearance of distance and depth in a detector's registration world is produced by bounded causal availability.
\end{proposition}

\begin{proof}
Events at different spatial relations to a detector enter its causal past through different null and timelike relations. Their information becomes locally available according to those relations, not all at once. The detector therefore receives a structured ordering of near and far, earlier and later, occluded and available, rather than an undifferentiated total state. By \cref{def:presented-depth}, this bounded ordering constitutes presented depth.
\end{proof}

\begin{remark}
The claim is not that geometric distance is fictitious. The claim is that perceived or recorded depth is never a direct apprehension of an extended universal present. It is the detector-side organization of causally available information.
\end{remark}

\begin{proposition}[The causal bound structures the registration world]
\label{prop:world-structure}
The invariant causal bound structures which relations can be jointly presented at one detector event and therefore structures the detector's registration world.
\end{proposition}

\begin{proof}
By \cref{prop:finite-access}, only the causal past contributes information to the detector event. By \cref{prop:depth}, the differential causal availability of events supplies presented depth and ordering. The invariant bound \(c\) determines the null structure delimiting that causal availability. Therefore the bound structures which relations can be jointly presented and how they appear as an ordered world.
\end{proof}

This supplies the stronger TLM reading:

\begin{quote}
The information speed limit does not merely restrict messages moving through a pre-given world. It is part of the lawful architecture by which an extended world is presented locally at all.
\end{quote}

\section{Why the Macroscopic Present Appears Stable}

The absence of a universal present does not make the macroscopic world chaotic. Large material systems preserve stable relational structure across enormous numbers of registrations.

\begin{definition}[Macroscopic registration stability]
\label{def:macro-stability}
Macroscopic registration stability is the persistence of mutually consistent relational patterns across repeated registrations of a large system.
\end{definition}

\begin{proposition}[Stable relational presentation]
\label{prop:stable-presentation}
A planet, table, building, or organism appears continuously present because immense numbers of micro-local presents remain lawfully coordinated into a highly stable relational pattern across repeated registrations.
\end{proposition}

\begin{proof}
A macroscopic body participates in vast numbers of interactions and registrations involving electromagnetic, gravitational, thermal, mechanical, and other relations. Across ordinary intervals, these registrations preserve approximately the same large-scale position, shape, momentum, composition, and relation to neighboring systems. The detector does not require an infinitely precise universal state of the body. It requires successive local records whose lawful relations remain stable. By \cref{def:macro-stability}, this produces macroscopic registration stability and the appearance of continuous presence.
\end{proof}

\begin{remark}
A planet has a stable orbital relation to its star. That stability is real. The planet does not possess one monolithic present. It consists of an immense plurality of micro-local presents whose lawful relations remain stable enough to register as one persistent planet. The reality lies in that durable coordination, not in a universe-wide instantaneous material tableau.
\end{remark}

\begin{corollary}[Persistence without universal simultaneity]
\label{cor:persistence}
Macroscopic persistence does not require a universal present.
\end{corollary}

\begin{proof}
By \cref{prop:stable-presentation}, persistence is presented through stable lawful relations across local registrations. None of those registrations requires a unique observer-independent assignment of distant simultaneity. Therefore macroscopic persistence does not require a universal present.
\end{proof}

\section{Quantum Outcomes and the Thinness of the Present}

At small scales, standard quantum theory does not generally assign a complete set of simultaneously definite classical properties prior to measurement. It assigns a state and probability structure from which outcome statistics are calculated~\cite{Dirac1958,Sakurai2017}. The detector records a definite result at registration.

The present paper does not attempt a complete interpretation of quantum mechanics. It uses only the limited point required for the registration claim: the formalism supplies probabilities for possible registered outcomes, while the detector event supplies the definite local record.

\begin{definition}[Lawful outcome]
\label{def:lawful-outcome}
A lawful outcome is a determinate detector registration whose admissibility and probability are fixed by the applicable quantum state, dynamics, boundary conditions, conservation laws, and detector interaction.
\end{definition}

\begin{proposition}[Probability-to-registration structure]
\label{prop:probability-registration}
At the detector scale, quantum theory relates a lawful probability structure to a definite local registration.
\end{proposition}

\begin{proof}
For a quantum state \(\lvert\psi\rangle\) and a measurement represented by projectors \(P_i\), the Born rule assigns outcome probabilities
\begin{equation}
p_i=\langle\psi|P_i|\psi\rangle.
\end{equation}
The probabilities describe the lawful distribution of possible registered outcomes. A completed detector interaction yields a definite record corresponding to one outcome in the measurement basis. Therefore the formal structure relates probability-governed admissibility to local registration.
\end{proof}

\begin{proposition}[No infinitely thick classical present]
\label{prop:no-thick-present}
The locally registered present need not contain an infinitely precise classical state of every constituent of the macroscopic world.
\end{proposition}

\begin{proof}
By \cref{prop:probability-registration}, the small-scale description is probability-governed and registration-specific rather than a universal inventory of simultaneously definite classical properties. By \cref{prop:stable-presentation}, macroscopic stability arises from repeated consistency of large-scale relations. Therefore the detector's present need not contain an infinitely precise classical state of every constituent. It contains the determinate local registration and the stable relational structure supported by the available records.
\end{proof}

This is the sense in which the present was ``barely here.'' The phrase does not mean that the present is weak, unreal, or momentary in a merely poetic sense. It means that the physically definite present is thinner than the classical picture suggests. It is not a complete universal material state. It is one local registration event carrying a lawfully structured presentation of a world.

\section{Where Is the Past?}

The ordinary question ``Where did the past go?'' assumes that the past once existed as a universal present and later departed. Under the registration-local account, that premise is rejected.

\begin{definition}[Prior registration]
\label{def:prior-registration}
A prior registration is a registration event lying earlier than the detector's current event along an applicable timelike or causal ordering.
\end{definition}

\begin{definition}[Present trace]
\label{def:present-trace}
A present trace of a prior registration is a current physical record, memory state, structural alteration, document, image, causal consequence, or other presently instantiated condition whose content refers to or was produced by the prior registration.
\end{definition}

\begin{proposition}[The past was never universal]
\label{prop:past-not-universal}
The past was never a single universal present that later ceased to exist.
\end{proposition}

\begin{proof}
By \cref{cor:not-invariant}, no universal present is supplied even at the time assigned to an earlier event. What existed were local registration events and lawful relations among them. As later detector events occur, those earlier events stand in the detector's causal past. Therefore the past was never one universal present that later vanished.
\end{proof}

\begin{proposition}[Past availability through present traces]
\label{prop:past-traces}
A detector's physical access to the past consists of present traces instantiated at its current event.
\end{proposition}

\begin{proof}
A detector cannot occupy an earlier event while remaining at its current event. What it physically possesses now are memory states, records, incoming signals, material consequences, and other current conditions produced by or referring to earlier events. By \cref{def:present-trace}, these are present traces. Therefore physical access to the past consists of present traces instantiated now.
\end{proof}

\begin{remark}
Remembering an event from thirty years ago is not the retrieval of a vanished universal present from a storage location called ``the past.'' It is a current registration in a present nervous system whose lawful content refers to an earlier registration. The memory is present; its referent is earlier.
\end{remark}

\begin{corollary}[The past does not go anywhere]
\label{cor:past-nowhere}
The past does not ``go'' anywhere, because it was never an extended substance or universal state capable of relocation.
\end{corollary}

\begin{proof}
By \cref{prop:past-not-universal}, the past consists of prior local events and their lawful relations, not a universal substance. Events do not need to travel elsewhere when later events occur. Their consequences remain available through present traces under \cref{prop:past-traces}. Therefore the past does not go anywhere.
\end{proof}

\section{The Shared World Is Consistency, Not Identity}

If every detector registers a different present, why do observers agree that they inhabit one world? They agree because the registrations are lawfully transformable and causally consistent.

\begin{proposition}[Consistency without identical presents]
\label{prop:consistency}
Two detectors can inhabit one shared world without registering identical local presents.
\end{proposition}

\begin{proof}
Let detectors \(A\) and \(B\) occupy different events. By \cref{prop:detector-relative}, they register different local presents. Their records can nevertheless be compared at later interactions. Relativistic coordinate transformations, causal ordering, conservation laws, field equations, and quantum statistics constrain the relation between the records. When those relations are satisfied, the records are mutually consistent even though they are not identical. By \cref{def:shared-world}, such consistency constitutes a shared world.
\end{proof}

\begin{proposition}[Agreement is a later registration]
\label{prop:agreement}
Agreement between observers is itself a local registration event, not proof that the observers previously occupied one identical present.
\end{proposition}

\begin{proof}
For two observers to compare records, information from each record must become locally available at a later interaction or communication event. The comparison produces a new physical record at that event. Their agreement therefore occurs as another local registration. It demonstrates consistency between prior records; it does not establish that the prior observers occupied one identical universal present.
\end{proof}

\begin{corollary}[Shared reality as lawful coherence]
\label{cor:shared-reality}
Shared physical reality is the lawful coherence of local registrations.
\end{corollary}

\begin{proof}
By \cref{prop:consistency}, different local presents form one shared world when their relations are mutually consistent. By \cref{prop:agreement}, inter-observer agreement is itself produced through further registration. Therefore shared physical reality is constituted operationally by the lawful coherence of local registrations, not by access to one universal present.
\end{proof}

\section{Central Result}

The argument can now be compressed.

\begin{proposition}[Local-present theorem]
\label{prop:local-present-theorem}
The present is not a universal physical condition. It is unique at every micro-local site in the universe. A detector or observer integrates many micro-local presents into a higher-order registration presented as a world, and every other detector or observer integrates a different present.
\end{proposition}

\begin{proof}
Relativity supplies no unique observer-independent distant simultaneity relation (\cref{prop:no-global-simultaneity,prop:gr-locality}). Physical locality is instantiated site by site (\cref{prop:registration-locality,prop:micro-uniqueness}). A detector is an extended coordination of many such sites (\cref{prop:detector-integration}). Different detectors occupy different events, possess different causal pasts, and integrate different collections of micro-local presents (\cref{prop:distinct-events,prop:different-pasts,prop:detector-relative}). Their registrations remain parts of one shared world because their relations are constrained to be lawfully consistent (\cref{prop:consistency,cor:shared-reality}). Therefore the present is not a universal physical condition. It is unique at every micro-local site, while each observer or detector integrates many micro-local presents into a distinct higher-order world-registration.
\end{proof}

\begin{corollary}[No pre-registered universe-wide now]
\label{cor:no-pre-registered-now}
The universe does not first possess one fully formed material present and then distribute partial views of it to detectors.
\end{corollary}

\begin{proof}
A pre-registered universe-wide now would be the universal present excluded by \cref{prop:local-present-theorem}. What detectors receive are different local registration worlds constrained by lawful relations. Therefore no pre-registered universe-wide now is required.
\end{proof}

\begin{corollary}[The present was barely here]
\label{cor:barely-here}
The present is physically thinner than the classical universal-now picture: it is one local determinate registration, not a complete observer-independent state of the universe.
\end{corollary}

\begin{proof}
By \cref{prop:local-present-theorem}, the present belongs to one local registration. By \cref{prop:no-thick-present}, that registration need not contain an infinitely precise classical state of every constituent. Therefore the physically definite present is thinner than the classical universal-now picture.
\end{proof}

\section{What the Paper Does Not Claim}

The result does not claim that nothing exists outside a human mind. Registration is physical and non-anthropic. Detectors include ordinary material systems.

The result does not claim that observers create arbitrary worlds. The probability structure, causal structure, conservation laws, field equations, detector couplings, and boundary conditions constrain every admissible registration.

The result does not claim that macroscopic bodies lack stability, location, momentum, or durable relations. It claims that these are presented through lawfully consistent registration patterns rather than through one observer-independent universal present.

The result does not deny coordinate descriptions of extended spacelike hypersurfaces. It denies that one chosen hypersurface is thereby promoted into the uniquely real physical present of the universe.

The result does not modify special relativity, general relativity, quantum mechanics, or quantum field theory. It changes only the ontology attached to their observer-relative, local, causal, and probability-governed structure.

\section{Discussion}

The universal-present picture survives in ordinary thought because the local world is extraordinarily coherent. Nearby registrations arrive in rapid succession. Macroscopic relations remain stable. Memories and records preserve prior events. Communication allows observers to compare results. The practical effect is a nearly seamless presentation of an extended world.

That seamlessness invites an ontological inflation. The observer mistakes the coherence of the registration world for proof that the entire universe possesses one simultaneous material state matching the observer's present. Relativity removes the simultaneity. Detector locality removes the global access. Quantum probability removes the requirement that every constituent possess a fully classical set of simultaneously registered properties. What remains is not nothing. What remains is lawful registration.

The information speed limit is central to this presentation. Because causal information is bounded, no detector receives all spatially separated events together. The detector receives a structured causal past. Distance appears through differential availability. Depth appears through ordered registration. Persistence appears through the recurrence of stable relational patterns. Shared history appears through present records whose causal relations remain consistent.

Under this reading, the information speed limit participates in the architecture of appearance. It is not merely a traffic law imposed on objects already moving through a fully presented universe. It is the structural limit under which a detector is presented with any extended world.

This reframes both past and present.

The past is not a discarded region into which the universal present has moved. It is the set of prior local registrations and lawful events available now only through present traces and causal consequences.

The present is not a total material inventory possessed by the universe. It is the local registration event at which a detector receives one lawfully structured world.

The future is not a universally existing later tableau waiting to become present. Relative to the detector's current event, it is the set of later admissible registrations not yet locally instantiated.

The resulting architecture is austere:

\begin{enumerate}
\item lawful structure constrains what can register;
\item causal structure constrains what can become locally available;
\item registration supplies the determinate local present;
\item repeated consistency supplies the appearance of a stable world;
\item comparison among registrations supplies shared history.
\end{enumerate}

The universe-wide present is not needed at any step.

\section{Conclusion}

The present was never a universal physical condition. Special relativity withholds a unique frame-independent simultaneity relation for spatially separated events. General relativity preserves local causal structure without restoring a unique global now. Detector physics supplies only local registration. Quantum formalism supplies lawful probabilities for possible outcomes and definite records at registration, not an infinitely precise classical inventory of the entire universe.

Accordingly, a present exists uniquely at every micro-local site. It belongs to that indivisible causal location, not to the universe as a whole. A detector or observer is an extended system that integrates many micro-local presents into one higher-order registration. Every other detector occupies a different collection of sites, possesses a different causal past, and integrates a different present.

The different presents do not fragment reality. Their relations remain constrained by the laws of special relativity, general relativity, quantum mechanics, conservation, and causal structure. What appears to be a shared present is therefore not one identical moment distributed across observers. It is a network of different local registrations whose relations remain mutually consistent.

Macroscopic objects appear continuously present because immense numbers of micro-local presents remain lawfully coordinated across repeated registrations. A planet retains a durable relation to its star; a table retains its form; a record retains information about a prior event. These are real lawful continuities. They do not require a universal present.

The past did not begin as one universal present and then disappear. It consisted of local events whose consequences are available through present traces. The present likewise is not a universe-wide state waiting to be observed. It is the act of local registration itself.

The central result is therefore:

\begin{quote}
\textbf{The present was never a universal physical condition. It is unique at every micro-local site in the universe. An observer or detector integrates many micro-local presents into a higher-order registration presented as a world, and every other observer or detector integrates a different present.}
\end{quote}

The world is not first universally present and then observed. Local presence begins at every indivisible site, and the observer's world is the lawful integration of those micro-local presents into one coherent registration.

\begin{thebibliography}{99}

\bibitem{Einstein1905}
A.~Einstein,
``On the Electrodynamics of Moving Bodies,''
\emph{Annalen der Physik} \textbf{17}, 891--921 (1905).

\bibitem{Rindler2006}
W.~Rindler,
\emph{Relativity: Special, General, and Cosmological},
2nd ed.,
Oxford University Press (2006).

\bibitem{Wald1984}
R.~M.~Wald,
\emph{General Relativity},
University of Chicago Press (1984).

\bibitem{MTW1973}
C.~W.~Misner, K.~S.~Thorne, and J.~A.~Wheeler,
\emph{Gravitation},
W.~H.~Freeman (1973).

\bibitem{Malament2012}
D.~B.~Malament,
\emph{Topics in the Foundations of General Relativity and Newtonian Gravitation Theory},
University of Chicago Press (2012).

\bibitem{Dirac1958}
P.~A.~M.~Dirac,
\emph{The Principles of Quantum Mechanics},
4th ed.,
Oxford University Press (1958).

\bibitem{Sakurai2017}
J.~J.~Sakurai and J.~Napolitano,
\emph{Modern Quantum Mechanics},
2nd ed.,
Cambridge University Press (2017).

\bibitem{Bedrock}
J.~C.~W.~McKinley,
\emph{The Timeless Light Model: Bedrock Statement},
Zenodo (2026).
\doi{10.5281/zenodo.19167403}.

\bibitem{Timebound}
J.~C.~W.~McKinley,
\emph{Timebound Does Not Mean Traveler:
A No-Go on Deriving Bead-Path Ontology from Quantum Admissibility},
Zenodo (2026).
\doi{10.5281/zenodo.20114078}.

\bibitem{RegistrationBound}
J.~C.~W.~McKinley,
\emph{\(c\) Is a Registration Bound, Not a Traveler's Speed:
An Observer-Side Reading of the Information Speed Limit},
Zenodo (2026).
\doi{10.5281/zenodo.20175517}.

\bibitem{ObserverPerceivability}
J.~C.~W.~McKinley,
\emph{The Information Speed Limit as the Bound on Observer-Perceivability
in the Timeless Light Model},
Zenodo (2026).
DOI placeholder.

\end{thebibliography}

\end{document}
